%% kapa: exam=1 round=37 year=2026 period=2 form=A subject=accounting subject_ko=회계학 seq=8 number=48
%% source: Q-Net 감정평가사 기출문제 (cst00309) artlSeq=5258924 file=raw/1008/5258924/2026년 제37회 감정평가사 1차 2교시.pdf
%% conversion: kordoc 4.24.51 -> tools/gen1.py (auto); review=proofread-sonnet; flags=none
\begin{question}{48}
(주)감평은 20×1년 초 사용목적의 자원재생설비(내용연수 5년, 잔존가치 ￦0, 정액법 상각, 원가모형 적용)를 ￦18,000에 건설하였다. 동 설비는 내용연수가 종료되면 원상 복구할 의무가 있으며, 복구시점에 ￦5,000이 지출될 것으로 추정되었다.
(주)감평은 20×1년 말 동 설비의 복구비용 지출금액을 ￦8,000으로 변경하였다. 복구비용은 충당부채의 인식요건을 충족하며, 복구비용의 현재가치 계산에 적용할 할인율은 연 6 \%이며 변동되지 않는다. 동 설비와 관련하여 20×2년도 포괄손익계산서에 인식할 총비용은? (단, 단일금액 ￦1의 현가계수(5년, 6 \%)는 0.7473, 현가계수(4년, 6 \%)는 0.7921이며, 화폐금액은 소수점 첫째자리에서 반올림하고 단수차이로 인한 오차가 있으면 가장 근사치를 선택한다.) \par
\begin{choices}
  \item ￦5,322
  \item ￦4,942
  \item ￦4,867
  \item ￦4,584
  \item ￦4,347
\end{choices}
\end{question}
\answer{1}
